\section*{Глава 3. Разработка веб-приложения}
\addcontentsline{toc}{section}{Глава 3. Разработка веб-приложения}
\stepcounter{section}

Разработанная система ориентирована на хранение, систематизацию и выдачу учебных материалов с учетом роли пользователя. В ходе разработки были реализованы базовая структура проекта, аутентификация, авторизация, механизм проверки прав доступа, CRUD-интерфейс для основных сущностей, загрузка и скачивание файлов, а также формирование статистики по материалам.

\subsection{Разработка базовой структуры приложения и шаблонов страниц}

Была сформирована базовая структура проекта, обеспечивающая разделение логики приложения, описания моделей данных, маршрутов, шаблонов страниц и конфигурации. Такой подход упрощает сопровождение программного кода и делает проект более понятным.

Структура приложения представлена следующими основными компонентами:

\begin{itemize}[leftmargin=1.25cm]
    \item файл \texttt{run.py} --- точка запуска веб-приложения;
    \item файл \texttt{config.py} --- конфигурация системы, включая путь к базе данных, папке загрузок и список допустимых расширений файлов;
    \item модуль \texttt{app/\_\_init\_\_.py} --- инициализация Flask-приложения, объекта базы данных, менеджера авторизации и вспомогательных шаблонных фильтров;
    \item модуль \texttt{app/models.py} --- описание сущностей предметной области и связей между ними;
    \item модуль \texttt{app/routes.py} --- обработка HTTP-запросов и реализация пользовательских сценариев;
    \item модуль \texttt{app/decorators.py} --- декораторы для проверки ролей пользователей;
    \item каталог \texttt{app/templates/} --- HTML-шаблоны страниц;
    \item каталог \texttt{app/uploads/} --- сохранение загруженных учебных файлов.
\end{itemize}

В файле конфигурации были заданы ключевые параметры приложения: секретный ключ, адрес базы данных SQLite, путь к каталогу загрузки файлов, максимальный размер загружаемого файла и набор разрешенных расширений. Это обеспечивает централизованное управление настройками и упрощает перенос приложения в другую среду.

Базовые шаблоны страниц построены по единому принципу: каждая страница наследуется от общего базового шаблона, содержащего навигационное меню, блок вывода сообщений системы и общую структуру пользовательского интерфейса. В результате был сформирован единый стиль страниц и обеспечено удобство навигации по разделам приложения.

В интерфейсе реализованы следующие основные страницы:

\begin{center}
\small\begin{tabularx}{0.96\textwidth}{>{\raggedright\arraybackslash}p{3.8cm} >{\raggedright\arraybackslash}X >{\raggedright\arraybackslash}p{4.3cm}}
\toprule
\textbf{Страница} & \textbf{Назначение} & \textbf{Кому доступна} \\
\midrule
Главная страница & Краткая информация о системе и последние добавленные материалы & Всем авторизованным пользователям \\
Страница входа & Ввод логина и пароля для начала работы в системе & Неавторизованным пользователям \\
Пользователи & Просмотр, создание, редактирование и удаление учетных записей & Администратору \\
Категории & Просмотр и управление категориями учебных материалов & Администратору и преподавателю \\
Материалы & Просмотр списка материалов, фильтрация и переход к карточке материала & Всем авторизованным пользователям \\
Карточка материала & Просмотр описания материала и прикрепленных файлов & Всем авторизованным пользователям \\
Статистика & Агрегированные сведения по материалам, пользователям и файлам & Администратору \\
\bottomrule
\end{tabularx}
\end{center}

Таким образом, была создана базовая структура Flask-приложения и набор шаблонов, необходимых для дальнейшей реализации предметной логики.

\subsection{Реализация аутентификации и авторизации пользователей}

Следующим этапом была реализована аутентификация пользователей. Для этого использована библиотека Flask-Login, позволяющая организовать сеансы работы пользователя, хранить состояние входа и ограничивать доступ к защищенным страницам.

В приложении реализована форма входа, содержащая поля логина и пароля. После отправки формы выполняется поиск пользователя по логину и проверка корректности пароля с использованием хэширования. В случае успешной проверки создается пользовательская сессия, после чего пользователь перенаправляется на главную страницу приложения. При ошибочном вводе данных выводится соответствующее сообщение.

Для хранения паролей используется не открытый текст, а хэш, вычисляемый средствами библиотеки Werkzeug. Это позволяет повысить безопасность приложения и соответствует типовой практике разработки веб-систем.

Ниже приведен фрагмент кода, реализующий проверку пароля в модели пользователя:

\begin{lstlisting}
def set_password(self, password: str) -> None:
    self.password_hash = generate_password_hash(password)

def check_password(self, password: str) -> bool:
    return check_password_hash(self.password_hash, password)
\end{lstlisting}

Для работы Flask-Login в проекте настроен загрузчик пользователя по идентификатору, что обеспечивает восстановление сведений о текущем пользователе из пользовательской сессии:

\begin{lstlisting}
@login_manager.user_loader
def load_user(user_id):
    return User.query.get(int(user_id))
\end{lstlisting}

После успешного входа пользователь получает доступ только к тем страницам и действиям, которые разрешены его ролью. При попытке обращения к защищенным маршрутам без авторизации система выдает сообщение о необходимости входа в систему.

Таким образом, в приложении была реализована полноценная аутентификация пользователей и базовая авторизация на уровне пользовательской сессии.

\subsection{Реализация системы проверки прав доступа в соответствии с ролью пользователя}

В разработанном приложении реализованы три роли:

\begin{itemize}[leftmargin=1.25cm]
    \item \texttt{admin} --- администратор;
    \item \texttt{teacher} --- преподаватель;
    \item \texttt{student} --- студент.
\end{itemize}

Роль пользователя хранится в отдельной сущности \texttt{Role}, а в сущности \texttt{User} предусмотрена связь с соответствующей записью роли. Для проверки ролей в модели пользователя реализован метод \texttt{has\_role()}, возвращающий истинное значение, если роль пользователя входит в заданный набор.

Для ограничения доступа к маршрутам был создан собственный декоратор \texttt{roles\_required}. Он используется совместно с декоратором \texttt{@login\_required} и проверяет, обладает ли текущий пользователь допустимой ролью. Если пользователь не авторизован, возвращается код 401, а если его роль не соответствует требуемой, возвращается код 403.

\begin{lstlisting}
def roles_required(*role_names):
    def decorator(view_func):
        @wraps(view_func)
        def wrapped(*args, **kwargs):
            if not current_user.is_authenticated:
                abort(401)
            if not current_user.has_role(*role_names):
                abort(403)
            return view_func(*args, **kwargs)
        return wrapped
    return decorator
\end{lstlisting}

Кроме проверки роли на уровне маршрута, для операций над материалами реализована дополнительная проверка принадлежности материала пользователю. Это необходимо для того, чтобы преподаватель мог редактировать и удалять только собственные материалы, тогда как администратор может управлять всеми материалами.

Следовательно, в приложении реализовано ролевое разграничение доступа, полностью соответствующее требованиям.

\subsection{Реализация CRUD-интерфейса для пользователей, категорий и материалов}

В приложении реализованы операции создания, просмотра, редактирования и удаления для трех основных сущностей: пользователей, категорий и материалов. Все операции доступны через веб-интерфейс и соответствующие маршруты Flask.

\subsubsection*{CRUD для пользователей}

Управление пользователями доступно только администратору. Он может просматривать список учетных записей, создавать новых пользователей, изменять логин, ФИО, пароль и роль, а также удалять пользователей. Для защиты системы предусмотрено ограничение, запрещающее администратору удалить собственную учетную запись во время работы в системе.

\subsubsection*{CRUD для категорий}

Категории используются для систематизации учебных материалов. Администратор может создавать новые категории, изменять их название и описание, а также удалять категории. При этом реализована проверка, запрещающая удаление категории, если с ней уже связаны материалы. Такое ограничение поддерживает целостность данных.

\subsubsection*{CRUD для материалов}

Материал является основной сущностью приложения. При создании материала указываются его название, описание и категория. Автором материала автоматически становится текущий авторизованный пользователь. После создания материал отображается в общем списке и доступен для просмотра пользователям с соответствующими правами.

При редактировании материала допускается изменение названия, описания и категории, а также добавление новых файлов. В режиме редактирования отображается список уже загруженных файлов, для которых реализованы отдельные действия по скачиванию и удалению.

Основные CRUD-маршруты приложения представлены в таблице.

\begin{center}
\footnotesize\begin{tabularx}{0.96\textwidth}{>{\raggedright\arraybackslash}p{6.0cm} >{\raggedright\arraybackslash}X >{\raggedright\arraybackslash}p{1.6cm}}
\toprule
\textbf{Маршрут} & \textbf{Назначение} & \textbf{Методы} \\
\midrule
/users & Список пользователей & GET \\
/users/create & Создание пользователя & GET, POST \\
/users/{id}/edit & Редактирование пользователя & GET, POST \\
/users/{id}/delete & Удаление пользователя & POST \\
/categories & Список категорий & GET \\
/categories/create & Создание категории & GET, POST \\
/categories/{id}/edit & Редактирование категории & GET, POST \\
/categories/{id}/delete & Удаление категории & POST \\
/materials & Список материалов и фильтрация & GET \\
/materials/create & Создание материала & GET, POST \\
/materials/{id} & Карточка материала & GET \\
/materials/{id}/edit & Редактирование материала & GET, POST \\
/materials/{id}/delete & Удаление материала & POST \\
\bottomrule
\end{tabularx}
\end{center}

Таким образом, в системе реализован полноценный CRUD-интерфейс для основных сущностей предметной области.

\subsection{Реализация загрузки, скачивания и удаления файлов учебных материалов}

Одной из ключевых функций приложения является работа с учебными файлами. Для этого в проекте реализована отдельная сущность \texttt{StoredFile}, содержащая оригинальное имя файла, внутреннее имя для хранения, путь к файлу на диске, дату загрузки и связь с материалом.

При загрузке файла выполняется проверка расширения. Разрешены только типы файлов, соответствующие учебным материалам: PDF, DOC, DOCX, TXT, PPT, PPTX, XLS, XLSX, PNG, JPG, JPEG. После проверки файл сохраняется в каталог \texttt{app/uploads}. Для предотвращения конфликтов имен и перезаписи существующих файлов используется генерация уникального имени на основе UUID.

Фрагмент кода сохранения файлов имеет следующий вид:

\begin{lstlisting}
original_name = secure_filename(file.filename)
extension = original_name.rsplit('.', 1)[1].lower()
stored_name = f"{uuid.uuid4().hex}.{extension}"
file_path = os.path.join(upload_folder, stored_name)
file.save(file_path)
\end{lstlisting}

Скачивание файла реализовано отдельным маршрутом. При обращении к нему сервер находит файл по идентификатору и отправляет его пользователю как вложение, сохраняя исходное имя файла для скачивания.

Удаление файла также выполняется отдельным маршрутом. При этом сначала удаляется физический файл с диска, затем удаляется соответствующая запись из базы данных. Возможность удаления предоставляется только администратору или преподавателю, имеющему право управления данным материалом.

Реализованный механизм загрузки, скачивания и удаления файлов позволяет полностью закрыть одно из обязательных требований, связанное с обработкой файловых данных.

\subsection{Реализация агрегирования данных и формирования статистики по материалам}

Для выполнения требования, связанного с обработкой и агрегированием данных, в приложении была реализована отдельная страница статистики. Доступ к ней предоставляется только администратору.

На странице статистики формируются следующие агрегированные показатели:

\begin{itemize}[leftmargin=1.25cm]
    \item общее количество материалов в системе;
    \item общее количество загруженных файлов;
    \item количество материалов по категориям;
    \item количество материалов по авторам;
    \item количество пользователей по ролям.
\end{itemize}

Для формирования статистики используются агрегатные функции SQLAlchemy, в частности \texttt{COUNT()}, а также операции группировки по категориям, авторам и ролям. Пример формирования статистики по категориям представлен ниже:

\begin{lstlisting}
materials_by_category = (
    db.session.query(Category.name, func.count(Material.id))
    .outerjoin(Material, Material.category_id == Category.id)
    .group_by(Category.id)
    .order_by(Category.name)
    .all()
)
\end{lstlisting}

Полученные значения передаются в шаблон страницы статистики и выводятся в виде наглядных таблиц. Это позволяет администратору оценивать наполнение архива, распределение материалов по категориям и активность пользователей.

\subsection{Тестирование основных пользовательских сценариев и проверка прав доступа}

После завершения реализации основных модулей приложения было выполнено функциональное тестирование пользовательских сценариев. Целью тестирования являлась проверка корректности работы разработанного веб-приложения в условиях, приближенных к его реальному использованию, а также подтверждение выполнения требований курсового проекта, связанных с аутентификацией, авторизацией, разграничением прав доступа, CRUD-операциями, обработкой файлов и формированием статистики.

Тестирование проводилось методом последовательной ручной проверки ключевых пользовательских сценариев. В качестве тестовых ролей использовались три учетные записи: администратор, преподаватель и студент. Проверка включала как позитивные сценарии, при которых пользователь выполняет допустимые действия, так и негативные сценарии, при которых система должна запретить доступ к функции или странице. Основное внимание уделялось корректности маршрутизации, сохранению данных в базе, изменению интерфейса в зависимости от роли пользователя, а также работе операций загрузки, скачивания и удаления файлов.

Для проверки были подготовлены тестовые данные: несколько пользователей с различными ролями, набор категорий учебных материалов, несколько записей материалов и прикрепленные к ним файлы разных типов. После каждого действия проверялось не только визуальное изменение данных на странице, но и фактическое наличие или отсутствие соответствующих записей в базе данных и файловом хранилище приложения.

Результаты тестирования основных сценариев приведены ниже.

\begin{longtable}{>{\centering\arraybackslash}p{0.9cm} >{\raggedright\arraybackslash}p{5.5cm} >{\raggedright\arraybackslash}p{4.2cm} >{\raggedright\arraybackslash}p{4.0cm}}
\toprule
\textbf{№} & \textbf{Проверяемый сценарий} & \textbf{Ожидаемый результат} & \textbf{Фактический результат} \\
\midrule
\endfirsthead
\toprule
\textbf{№} & \textbf{Проверяемый сценарий} & \textbf{Ожидаемый результат} & \textbf{Фактический результат} \\
\midrule
\endhead
1 & Открытие защищенной страницы без авторизации & Система перенаправляет пользователя на страницу входа & Выполнено, доступ без входа не предоставляется \\
2 & Вход администратора по корректным учетным данным & Открывается главная страница приложения, пользовательская сессия создается & Выполнено, вход успешен \\
3 & Выход из системы & Сессия пользователя завершается, защищенные страницы снова недоступны & Выполнено, повторный доступ требует авторизации \\
4 & Создание нового пользователя администратором & В базе появляется новая учетная запись с выбранной ролью & Выполнено, пользователь создается и отображается в списке \\
5 & Редактирование данных пользователя администратором & Измененные значения сохраняются и отображаются в интерфейсе & Выполнено, обновленные данные отображаются корректно \\
6 & Удаление пользователя администратором & Учетная запись удаляется из базы и исчезает из списка & Выполнено, запись удаляется корректно \\
7 & Создание категории администратором & Новая категория сохраняется и доступна при создании материалов & Выполнено, категория появляется в списке \\
8 & Попытка студента открыть страницу категорий & Доступ запрещается, страница недоступна студенту & Выполнено, доступ студенту ограничен \\
9 & Создание нового материала преподавателем & Материал сохраняется с указанием автора и категории & Выполнено, материал добавляется в систему \\
10 & Редактирование материала его автором & Изменения сохраняются и отображаются на странице просмотра материала & Выполнено, материал обновляется корректно \\
11 & Попытка студента изменить существующий материал & Система запрещает выполнение операции & Выполнено, права студента ограничены \\
12 & Загрузка файла при создании или редактировании материала & Файл сохраняется на сервере и отображается в карточке материала & Выполнено, файл доступен в списке вложений \\
13 & Скачивание прикрепленного файла авторизованным пользователем & Браузер начинает загрузку выбранного файла & Выполнено, файл скачивается корректно \\
14 & Удаление файла из карточки материала пользователем с соответствующими правами & Файл удаляется из хранилища и перестает отображаться в интерфейсе & Выполнено, удаление работает корректно \\
15 & Открытие страницы статистики администратором & Отображаются агрегированные показатели по материалам, файлам и пользователям & Выполнено, статистика формируется корректно \\
16 & Попытка преподавателя открыть страницу статистики & Доступ запрещается, страница статистики недоступна преподавателю & Выполнено, доступ ограничен \\
17 & Попытка студента открыть страницу статистики & Доступ запрещается, страница статистики недоступна студенту & Выполнено, доступ ограничен \\
18 & Проверка отображения автора материала в интерфейсе & Для каждого пользователя отображается формат ``ФИО (логин)'' & Выполнено, отображение авторов единообразно \\
19 & Проверка отображения времени создания записей & Дата и время выводятся с учетом часового пояса МСК (UTC+3) & Выполнено, время отображается корректно \\
20 & Проверка агрегированных данных после добавления новых материалов и файлов & Количественные показатели изменяются в соответствии с новыми данными & Выполнено, статистика обновляется корректно \\
\bottomrule
\end{longtable}

По результатам проведенного тестирования можно сделать вывод, что основные пользовательские сценарии выполняются корректно. Проверка показала, что приложение устойчиво отрабатывает как обычные операции пользователя, так и ограничения, связанные с ролью учетной записи. Особенно важным результатом является подтверждение того, что студент не имеет доступа к административным функциям, преподаватель не может просматривать статистику, а изменение и удаление материалов и файлов доступно только пользователям, обладающим соответствующими правами.

Также в процессе тестирования были выявлены и устранены отдельные недочеты интерфейса и серверной логики. К ним относились вопросы отображения времени, единообразного показа имени пользователя в формате ``ФИО (логин)'', а также корректная обработка удаления уже загруженных файлов со страницы редактирования материала. После внесения исправлений указанные сценарии были протестированы повторно и подтвердили корректность реализации.